
%\addcontentsline{toc}{section}{Б.\,Дидактичні матеріали до практичної роботи №2}

% \addtocontents{toc}{\protect\setcounter{tocdepth}{-1}}


\section {\centering Дидактичні матеріали до практичної роботи №2}

% Налаштування для коду Python
\lstset{
    language=Python,
    basicstyle=\ttfamily\small,
    commentstyle=\color{gray},
    keywordstyle=\color{blue},
    stringstyle=\color{red},
    numbers=left,
    numberstyle=\tiny\color{gray},
    stepnumber=1,
    numbersep=5pt,
    backgroundcolor=\color{white},
    showspaces=false,
    showstringspaces=false,
    showtabs=false,
    frame=single,
    rulecolor=\color{black},
    tabsize=4,
    captionpos=b,
    breaklines=true,
    breakatwhitespace=false,
    escapeinside={(*@}{@*)},
    inputencoding=utf8,
    extendedchars=true,
    literate={а}{{\cyra}}1 {б}{{\cyrb}}1 {в}{{\cyrv}}1 {г}{{\cyrg}}1 {д}{{\cyrd}}1
             {е}{{\cyre}}1 {є}{{\cyrie}}1 {ж}{{\cyrzh}}1 {з}{{\cyrz}}1 {и}{{\cyri}}1
             {і}{{\cyrii}}1 {ї}{{\cyryi}}1 {й}{{\cyrishrt}}1 {к}{{\cyrk}}1 {л}{{\cyrl}}1
             {м}{{\cyrm}}1 {н}{{\cyrn}}1 {о}{{\cyro}}1 {п}{{\cyrp}}1 {р}{{\cyrr}}1
             {с}{{\cyrs}}1 {т}{{\cyrt}}1 {у}{{\cyru}}1 {ф}{{\cyrf}}1 {х}{{\cyrh}}1
             {ц}{{\cyrc}}1 {ч}{{\cyrch}}1 {ш}{{\cyrsh}}1 {щ}{{\cyrshch}}1 {ь}{{\cyrsftsn}}1
             {ю}{{\cyryu}}1 {я}{{\cyrya}}1
}



\section*{\centering Академічний напрямок}

\subsection*{\centering ЛАБОРАТОРНИЙ ЖУРНАЛ\\Дослідження криптографічних алгоритмів}

\noindent
\textbf{ПІБ учня:} \rule{6cm}{0.5pt}\\
\textbf{Дата проведення:} \rule{5cm}{0.5pt}\\
\textbf{Обрані інструменти:} \rule{5cm}{0.5pt}

\subsubsection*{Експеримент 1. Порівняння хеш-функцій}

\noindent
\textbf{Вхідні дані:} \rule{10cm}{0.5pt}

\begin{table}[h]
\centering
\begin{tabular}{|l|p{5cm}|p{2cm}|p{1.5cm}|p{2.5cm}|}
\hline
\textbf{Алгоритм} & \textbf{Результат хешування} & \textbf{Довжина (біт)} & \textbf{Час виконання (мс)} & \textbf{Примітки} \\
\hline
MD5 & & 128 & & \\
\hline
SHA-1 & & 160 & & \\
\hline
SHA-256 & & 256 & & \\
\hline
\end{tabular}
\end{table}

\noindent
\textbf{Спостереження:} \rule{12cm}{0.5pt}\\
\rule{16cm}{0.5pt}\\
\rule{16cm}{0.5pt}\\
\rule{16cm}{0.5pt}

\subsubsection*{Експеримент 2. Лавинний ефект}

\noindent
\textbf{Оригінальний текст:} \rule{6cm}{0.5pt}\\
\textbf{Модифікований текст (зміна 1 символу):} \rule{6cm}{0.5pt}

\begin{table}[h]
\centering
\begin{tabular}{|l|p{3cm}|p{3cm}|c|}
\hline
\textbf{Алгоритм} & \textbf{Хеш оригіналу} & \textbf{Хеш модифікованого} & \textbf{Змінених біт (\%)} \\
\hline
MD5 & & & \\
\hline
SHA-256 & & & \\
\hline
\end{tabular}
\end{table}

\noindent
\textbf{Висновок про лавинний ефект:} \rule{10cm}{0.5pt}\\
\rule{15cm}{0.5pt}

\subsubsection*{Експеримент 3. Криптоаналіз (для достатнього рівня)}

\textbf{Тестові хеші для злому:}
\begin{enumerate}
    \item 5f4dcc3b5aa765d61d8327deb882cf99 (MD5)
    \item 25d55ad283aa400af464c76d713c07ad (MD5)
    \item 827ccb0eea8a706c4c34a16891f84e7b (MD5)
\end{enumerate}

\begin{table}[h]
\centering
\begin{tabular}{|c|l|l|c|}
\hline
\textbf{Хеш} & \textbf{Знайдений пароль} & \textbf{Метод злому} & \textbf{Час на злом} \\
\hline
1 & & & \\
\hline
2 & & & \\
\hline
3 & & & \\
\hline
\end{tabular}
\end{table}

\noindent
\textbf{Рекомендації щодо захисту паролів:} \rule{9cm}{0.5pt}\\
\rule{15cm}{0.5pt}

\subsubsection*{Експеримент 4. Частотний аналіз шифру Цезаря}

\noindent
\textbf{Зашифроване повідомлення:} \rule{10cm}{0.5pt}\\
\rule{15cm}{0.5pt}

\textbf{Частотна таблиця символів:}

\begin{table}[h]
\centering
\begin{tabular}{|c|c|c|c|}
\hline
\textbf{Символ} & \textbf{Кількість} & \textbf{Частота (\%)} & \textbf{Ймовірна заміна} \\
\hline
 & & & \\
\hline
 & & & \\
\hline
 & & & \\
\hline
\end{tabular}
\end{table}

\noindent
\textbf{Визначений зсув:} \rule{2cm}{0.5pt}\\
\textbf{Розшифроване повідомлення:} \rule{9cm}{0.5pt}\\
\rule{15cm}{0.5pt}

\subsubsection*{Загальні висновки з лабораторної роботи:}
\begin{enumerate}
    \item \textbf{Порівняння стійкості алгоритмів:} \rule{8cm}{0.5pt}
    \item \textbf{Практична небезпека слабких паролів:} \rule{7cm}{0.5pt}
    \item \textbf{Рекомендації для закладу освіти:} \rule{7.5cm}{0.5pt}
\end{enumerate}

\newpage
\section*{\centering Технічний напрямок}

\subsection*{Шаблон 1. Базове шифрування файлів (Python)}

\begin{lstlisting}[caption={Шаблон для симетричного шифрування файлів}]
----
шаблон для симетричного шифрування файлів
використовує бібліотеку cryptography
----

from cryptography.hazmat.primitives import hashes
from cryptography.hazmat.primitives.kdf.pbkdf2 import PBKDF2
from cryptography.hazmat.primitives.ciphers import Cipher, algorithms, modes
from cryptography.hazmat.backends import default_backend
import os

def generate_key_from_password(password: str, salt: bytes) -> bytes:
    """
    генерує криптографічний ключ з пароля
    
    параметри:
    - password: пароль користувача
    - salt: випадкова сіль для унікальності ключа
    
    повертає: 32-байтний ключ для AES-256
    """
    # TODO: використати PBKDF2 для генерації ключа
    # підказка: iterations=100000, length=32
    kdf = PBKDF2(
        algorithm=hashes.SHA256(),
        length=32,
        salt=salt,
        iterations=100000,
        backend=default_backend()
    )
    key = # ваш код тут
    return key

def encrypt_file(input_file: str, output_file: str, password: str):
    """
    шифрує файл за допомогою AES в режимі CBC
    
    параметри:
    - input_file: шлях до вихідного файлу
    - output_file: шлях для збереження зашифрованого файлу
    - password: пароль для шифрування
    """
    # генеруємо випадкову сіль та вектор ініціалізації
    salt = os.urandom(16)  # 16 байт випадкових даних
    iv = os.urandom(16)    # вектор ініціалізації для CBC
    
    # генеруємо ключ з пароля
    key = generate_key_from_password(password, salt)
    
    # читаємо дані з файлу
    with open(input_file, 'rb') as f:
        plaintext = f.read()
    
    # додаємо padding до даних (AES вимагає блоки по 16 байт)
    padding_length = 16 - (len(plaintext) % 16)
    padded_data = plaintext + bytes([padding_length] * padding_length)
    
    # TODO: створити об'єкт Cipher з AES та режимом CBC
    cipher = # ваш код тут
    encryptor = cipher.encryptor()
    
    # шифруємо дані
    ciphertext = # ваш код тут
    
    # зберігаємо зашифровані дані разом з salt та IV
    with open(output_file, 'wb') as f:
        f.write(salt)        # перші 16 байт - сіль
        f.write(iv)          # наступні 16 байт - IV
        f.write(ciphertext)  # решта - зашифровані дані
    
    print(f"файл зашифровано: {output_file}")

def decrypt_file(input_file: str, output_file: str, password: str):
    """
    розшифровує файл
    
    TODO: реалізувати зворотній процес до encrypt_file
    """
    # ваш код тут
    pass
\end{lstlisting}

\subsection*{Шаблон 2. Менеджер паролів (Python)}

\begin{lstlisting}[caption={Простий менеджер паролів з хешуванням}]
"""
простий менеджер паролів з хешуванням
демонстрація безпечного зберігання паролів
"""

import hashlib
import json
import getpass
from typing import Dict, Optional

class PasswordManager:
    """Kлас для безпечного управління паролями"""
    
    def __init__(self, master_password: str):
        """
        Iніціалізація менеджера з мастер-паролем
        
        master_password: головний пароль для доступу
        """
        # хешуємо мастер-пароль для перевірки
        self.master_hash = self._hash_password(master_password)
        self.storage: Dict[str, Dict] = {}
        self.authenticated = False
    
    def _hash_password(self, password: str, salt: str = "") -> str:
        """
        хешує пароль за допомогою SHA-256
        
        TODO: додати сіль для підвищення безпеки
        """
        # базова реалізація (потребує покращення)
        data = password + salt
        # ваш код тут: використати hashlib.sha256
        return hashlib.sha256(data.encode()).hexdigest()
    
    def authenticate(self, password: str) -> bool:
        """
        перевірка мастер-пароля
        
        повертає True якщо пароль правильний
        """
        # TODO: порівняти хеш введеного пароля з збереженим
        input_hash = # ваш код тут
        if input_hash == self.master_hash:
            self.authenticated = True
            return True
        return False
\end{lstlisting}

\subsection*{Шаблон 3. Secure Messenger (для високого рівня)}

\begin{lstlisting}[caption={Прототип захищеного месенджера}]
"""
прототип захищеного месенджера
використовує гібридне шифрування (RSA + AES)
"""

from cryptography.hazmat.primitives.asymmetric import rsa, padding
from cryptography.hazmat.primitives import serialization, hashes
from cryptography.hazmat.primitives.ciphers import Cipher, algorithms, modes
import os
import json
from datetime import datetime

class SecureMessenger:
    """клас для безпечного обміну повідомленнями"""
    
    def __init__(self, user_name: str):
        """
        ініціалізація месенджера для користувача
        
        user_name: ім'я користувача
        """
        self.user_name = user_name
        self.private_key = None
        self.public_key = None
        self.contacts = {}  # словник публічних ключів контактів
        self.message_history = []
        
        # генеруємо ключову пару RSA
        self._generate_key_pair()
    
    def _generate_key_pair(self):
        """
        генерує пару ключів RSA (2048 біт)
        
        TODO: реалізувати генерацію та збереження ключів
        """
        # генеруємо приватний ключ
        self.private_key = rsa.generate_private_key(
            public_exponent=65537,
            key_size=2048
        )
        
        # отримуємо публічний ключ
        self.public_key = # ваш код тут
        
        print(f"згенеровано ключі для {self.user_name}")
\end{lstlisting}

\newpage
\section*{\centering Управлінський напрямок}

\subsection*{КРИПТОГРАФІЧНА ПОЛІТИКА ОРГАНІЗАЦІЇ}
\textit{[Назва організації]}

\noindent
\textbf{Версія 1.0 | Дата затвердження:} \rule{4cm}{0.5pt}

\subsection*{1. ЗАГАЛЬНІ ПОЛОЖЕННЯ}

\subsubsection*{1.1. Мета документа}
Дана політика визначає стандарти та процедури використання криптографічних засобів захисту інформації в організації для забезпечення конфіденційності, цілісності та автентичності даних.

\subsubsection*{1.2. Сфера застосування}
Політика поширюється на всі інформаційні системи організації, включаючи:
\begin{itemize}
    \item Корпоративну електронну пошту
    \item Системи зберігання файлів та документів
    \item Бази даних з персональними та фінансовими даними
    \item Канали передачі даних між філіями
    \item Мобільні пристрої співробітників
\end{itemize}

\subsubsection*{1.3. Відповідальні особи}
\begin{itemize}
    \item \textbf{Власник політики:} Директор з інформаційної безпеки
    \item \textbf{Адміністратор криптографічних систем:} [ПІБ]
    \item \textbf{Контролюючий орган:} Комітет з інформаційної безпеки
\end{itemize}

\subsection*{2. КРИПТОГРАФІЧНІ СТАНДАРТИ}

\subsubsection*{2.1. Алгоритми шифрування}

\begin{table}[h]
\centering
\begin{tabular}{|l|l|l|l|}
\hline
\textbf{Призначення} & \textbf{Мінімальний} & \textbf{Рекомендований} & \textbf{Заборонені} \\
& \textbf{стандарт} & \textbf{стандарт} & \textbf{алгоритми} \\
\hline
Симетричне шифрування & AES-128 & AES-256 & DES, 3DES \\
\hline
Асиметричне шифрування & RSA-2048 & RSA-4096, ECC-384 & RSA < 2048 біт \\
\hline
Хешування & SHA-256 & SHA-512, SHA3-512 & MD5, SHA-1 \\
\hline
Цифровий підпис & RSA-2048 + SHA-256 & ECDSA з P-384 & DSA < 2048 біт \\
\hline
Обмін ключами & Diffie-Hellman 2048 & ECDH з P-384 & DH < 2048 біт \\
\hline
\end{tabular}
\end{table}

\subsubsection*{2.2. Генерація випадкових чисел}
Для генерації криптографічних ключів використовувати виключно апаратні або програмні генератори, сертифіковані за стандартом ДСТУ 4145-2002 або FIPS 140-2.

\subsection*{3. УПРАВЛІННЯ КЛЮЧАМИ}

\subsubsection*{3.1. Життєвий цикл ключів}

\begin{longtable}{|l|p{4cm}|l|l|}
\hline
\textbf{Етап} & \textbf{Процедура} & \textbf{Відповідальний} & \textbf{Термін} \\
\hline
\endfirsthead
\hline
\textbf{Етап} & \textbf{Процедура} & \textbf{Відповідальний} & \textbf{Термін} \\
\hline
\endhead
Генерація & Автоматизована через HSM або KMS & Адміністратор безпеки & За запитом \\
\hline
Розповсюдження & Захищені канали (TLS 1.2+) & Адміністратор безпеки & 24 години \\
\hline
Зберігання & Зашифроване сховище ключів & Система KMS & Постійно \\
\hline
Ротація & Планова заміна & Автоматизовано & Щорічно \\
\hline
Знищення & Безпечне видалення & Адміністратор безпеки & Після закінчення терміну \\
\hline
\end{longtable}

\subsubsection*{3.2. Захист ключів}
\begin{itemize}
    \item Майстер-ключі зберігаються в апаратному модулі безпеки (HSM)
    \item Робочі ключі шифруються майстер-ключами
    \item Резервні копії ключів зберігаються в окремому захищеному сховищі
    \item Доступ до ключів надається за принципом мінімальних привілеїв
\end{itemize}

\subsection*{4. СЦЕНАРІЇ ВИКОРИСТАННЯ}

\subsubsection*{4.1. Електронна пошта}
\begin{itemize}
    \item Обов'язкове використання S/MIME або PGP для листів з грифом ``Конфіденційно''
    \item Автоматичне шифрування вкладень понад 1 МБ
    \item Цифровий підпис для всіх зовнішніх комунікацій
\end{itemize}

\subsubsection*{4.2. Зберігання даних}
\begin{itemize}
    \item Повнодискове шифрування для ноутбуків (BitLocker, FileVault)
    \item Шифрування баз даних на рівні таблиць для критичних даних
    \item Шифрування резервних копій перед архівуванням
\end{itemize}

\subsubsection*{4.3. Передача даних}
\begin{itemize}
    \item TLS 1.2+ для всіх веб-сервісів
    \item VPN з IPSec для з'єднань між офісами
    \item SFTP/FTPS для передачі файлів
\end{itemize}

\subsection*{5. ПРОЦЕДУРИ РЕАГУВАННЯ}

\subsubsection*{5.1. Компрометація ключів}

Етапи реагування при підозрі на компрометацію:
\begin{enumerate}
    \item \textbf{Негайно} - блокування скомпрометованого ключа
    \item \textbf{1 година} - оповіщення всіх користувачів системи
    \item \textbf{4 години} - генерація та розповсюдження нових ключів
    \item \textbf{24 години} - повна заміна ключів у системі
    \item \textbf{48 годин} - аудит інцидента та звіт керівництву
\end{enumerate}

\subsubsection*{5.2. Втрата ключів}
\begin{itemize}
    \item Використання резервних копій з безпечного сховища
    \item Процедура аварійного розшифрування через ключі відновлення
    \item Документування інциденту в журналі безпеки
\end{itemize}

\subsection*{6. НАВЧАННЯ ТА АУДИТ}

\subsubsection*{6.1. Навчання персоналу}
\begin{itemize}
    \item Базовий курс криптографічної безпеки для всіх співробітників (щорічно)
    \item Спеціалізоване навчання для ІТ-персоналу (щопівроку)
    \item Тестування знань після навчання
\end{itemize}

\subsubsection*{6.2. Аудит}
\begin{itemize}
    \item Щоквартальна перевірка налаштувань криптографічних систем
    \item Щорічний зовнішній аудит криптографічної інфраструктури
    \item Моніторинг використання застарілих алгоритмів
\end{itemize}

\subsection*{7. КЛЮЧОВІ ПОКАЗНИКИ ЕФЕКТИВНОСТІ (KPI)}

\begin{table}[h]
\centering
\begin{tabular}{|p{6cm}|c|c|}
\hline
\textbf{Показник} & \textbf{Цільове значення} & \textbf{Частота вимірювання} \\
\hline
Відсоток зашифрованих критичних даних & $\geq$ 99\% & Щомісяця \\
\hline
Час реакції на компрометацію ключа & < 1 години & За інцидентом \\
\hline
Успішність ротації ключів & 100\% & Щорічно \\
\hline
Кількість інцидентів через слабку криптографію & 0 & Щоквартально \\
\hline
Проходження навчання персоналом & $\geq$ 95\% & Щорічно \\
\hline
\end{tabular}
\end{table}

\subsection*{8. НОРМАТИВНІ ПОСИЛАННЯ}

\begin{itemize}
    \item Закон України ``Про захист інформації в інформаційно-телекомунікаційних системах''
    \item ДСТУ 7624:2014 (Калина)
    \item ISO/IEC 27001:2013
    \item NIST SP 800-57 (Key Management)
\end{itemize}